% FILE: thesisheader.tex
% 
% by Scott MacGIBBON, Version 1.0 18/10/88.
% Adapted from header.tex on above date.
%
% Macros for chapter & section numbering, appendix numbering, 
% and headers and footers are below.
%
% Macros for lists are included.
%
\font\chapnumfont=amssmc40
%\font\chapnumfont=ambx10 scaled 1440
\font\bigbf=ambx10 scaled 1440
\font\bigrm=amr10 scaled 1440
\font\bigit=amsl10 scaled 1440
\font\medit=amsl10 scaled 1200
\font\medbf=ambx10 scaled 1200
\font\medrm=amr10 scaled 1200
\font\tenrm=amr10
\font\ninerm=amr9
\font\eightrm=amr8
\font\sevenrm=amr7
\font\ninett=amtt9
%
% Experimental setting of \parskip to try and increase distances between
% paragraphs.
%
\parskip=2.0pt plus 1.0pt
%
% and now try and increase spacing between lines...
%
\openup1\jot
%
% Setup macro for signalling that there is a chapter heading here.
%
\newif\iftitle
\def\titlepage{\global\titletrue}
\def\chapname{}
\def\negname{}
%
% Macros for headers and footers appears below.
%
\headline={\ifnum\pageno>0 
			\iftitle \global\titlefalse\nohead
			\else \normalhead \fi
	\else	\ifnum\pageno=0\nohead
		\else
			\iftitle \global\titlefalse\nohead
			\else \neghead \fi
	\fi	\fi}
\def\nohead{\hfill}
\def\normalhead{\tenrm\chapname \hfill Thesis on Brain Simulatio\rlap{n}}
\def\neghead{\tenrm\negname \hfill}
%
\footline={\ifnum\pageno=0\nofoot
	\else\normalfoot
	\fi}
\def\nofoot{\hfill}
\def\normalfoot{\hfill\tenrm\folio\hfill}
%%% The macros to generate sections and subsections (ssec's) etc
%%% are found below.
%
% First, set up the counting variables.
%
\newcount\partno	\global\partno=-1
\newcount\chapno	\global\chapno=-1
\newcount\secno		\global\secno=0
\newcount\ssecno	\global\ssecno=0
\newcount\sssecno	\global\sssecno=0
\newcount\ssssecno	\global\ssssecno=0
\newcount\appno		\global\appno=64
\newcount\sappno	\global\sappno=0
\newcount\ssappno	\global\ssappno=0
%
\newcount\lineno % for numbering lines in verbatim include files
%
% Setup output file for contents.
%
\newwrite\cntnts
\immediate\openout\cntnts=contents.tex
%
% The macro below, \contents, adds the table of contents.
%
\outer\def\contents {
	\immediate\closeout\cntnts
	\xdef\negname{\llap{C}ontents}
	\titlepage
	\openup-1\jot
	\noindent
	{\chapnumfont C}{\bigbf ontents}
	\bigskip
	\obeylines
	\input{contents}.tex
	\vfill\supereject
}
%
% The macro below, \ack, adds the acknowledgements
%
\outer\def\ack {
	\xdef\negname{\llap{A}cknowledgements}
	\titlepage
	\noindent
	{\chapnumfont A}{\bigbf cknowledgements}
	\bigskip
	\input{acknowledge}
	\vfill\supereject
}
%
% The chapter starting macros follow below.
%
% First, set up a boolean to indicate the first part, which doesn't
% require the \vfill\supereject.
%
\newif\iffirstpart
\global\firstparttrue
\outer\def\part #1 #2\par{
	\advance\partno by 1
	\iffirstpart\global\firstpartfalse
		\else\vfill\supereject\fi
	\immediate\write\cntnts{\string\smallskip}
	\immediate\write\cntnts{\string\noindent\string\bf\ Part\ \string\bf\the\partno\ #1\string\rm\string\dotfill\folio}
	\message{\the\partno}
	% Now write the chapter onto the page:
	\titlepage
	\vskip1cm
	\noindent
	{\bigbf \llap{P}art\ {\chapnumfont\the\partno:}#1\hfil}
	\bigskip
	{\sl\hfil #2}
}
%
%
\outer\def\chap #1\par{
	\global\secno=0		\global\ssecno=0
	\global\sssecno=0	\global\ssssecno=0
	\advance\chapno by 1
	\vfill\supereject
	\immediate\write\cntnts{\string\smallskip}
	\immediate\write\cntnts{\string\noindent\string\bf\the\chapno\string\bf\ #1\string\rm\string\dotfill\folio}
	\xdef\chapname{\llap{C}hapter\ \the\chapno. #1\unskip}
	\message{\the\chapno #1}
	% Now write the chapter onto the page:
	\titlepage
	\noindent
	{\bigbf \llap{C}hapter\ {\chapnumfont\the\chapno:} #1}
	\bigskip\noindent
}
%
% This is a section within the chapter macro.
%
\outer\def\sec #1\par{
	\global\ssecno=0
	\global\sssecno=0	\global\ssssecno=0
	\advance\secno by 1
	\immediate\write\cntnts{\string\indent\the\chapno.\the\secno\ #1\string\dotfill\folio}
	\message{\the\chapno.\the\secno}
	% Now write the section number onto the page:
	\vskip 0pt plus.2\vsize\penalty-250
	\medskip
	\vskip 0pt plus-0.2\vsize
	\noindent
	{\medbf\the\chapno.\the\secno. #1}
	\nobreak\smallskip\noindent
}
%
% This is a subsection macro within the chapter.
%
\outer\def\ssec #1\par{
	\global\sssecno=0	\global\ssssecno=0
	\advance\ssecno by 1
	\immediate\write\cntnts{\string\indent\string\indent\string\ninerm\the\chapno.\the\secno.\the\ssecno\ #1\string\dotfill\folio\string\tenrm}
	\message{\the\chapno.\the\secno.\the\ssecno}
	% Now write the section number onto the page:
	\vskip 0pt plus.1\vsize\penalty-250
	\smallskip
	\vskip 0pt plus-0.1\vsize
	\noindent
	{\bf\the\chapno.\the\secno.\the\ssecno. #1}
	\nobreak\smallskip\noindent
}
%
% This is a subsubsection macro within the chapter.
%
\outer\def\sssec #1\par{
	\global\ssssecno=0
	\advance\sssecno by 1
	\immediate\write\cntnts{\string\indent\string\indent\string\indent\string\eightrm\the\chapno.\the\secno.\the\ssecno.\the\sssecno\ #1\string\dotfill\folio\string\tenrm}
	% Now write the section number onto the page:
	\allowbreak\smallskip\noindent
	{\bf\the\chapno.\the\secno.\the\ssecno.\the\sssecno. #1:}
}
%
% This is a subsubsubsection macro within the chapter.
%
\outer\def\ssssec #1\par{
	\advance\ssssecno by 1
	\immediate\write\cntnts{\string\indent\string\indent\string\indent\string\indent\string\sevenrm\the\chapno.\the\secno.\the\ssecno.\the\sssecno.\the\ssssecno\ #1\string\dotfill\folio\string\tenrm}
	% Now write the section number onto the page:
	\allowbreak\smallskip\noindent
	{\sl\the\chapno.\the\secno.\the\ssecno.\the\sssecno.\the\ssssecno. #1:}
}
%
% The appendix starting macros follow below.
%
\outer\def\app #1 #2\par{
	\vfill\supereject
	\global\sappno=0	\global\ssappno=0
	\advance\appno by 1
	\xdef\appltr{#1}
	\nonfrenchspacing % resets spacing if references run first.
	\immediate\write\cntnts{\string\smallskip}
	\immediate\write\cntnts{\string\noindent\string\bf\ \appltr\string\bf\ #2\string\rm\string\dotfill\folio}
	\xdef\chapname{\llap{A}ppendix\ \char\appno. #2\unskip}
	\message{\appltr #2}
	% Now write the chapter onto the page:
	\titlepage
	\noindent
	{
		{\bigbf \llap{A}ppendix\ }
		{\chapnumfont \char\appno:}
		{\bigbf #2}
	}
	\bigskip\noindent
}
%
% This is a section within the appendix macro.
%
\outer\def\sapp #1\par{
	\global\ssappno=0
	\advance\sappno by 1
	\immediate\write\cntnts{\string\indent\ \appltr.\the\sappno\ #1\string\dotfill\folio}
	\message{\appltr.\the\sappno}
	% Now write the section number onto the page:
	\vskip 0pt plus.2\vsize\penalty-250
	\medskip
	\vskip 0pt plus-0.2\vsize
	\noindent
	{\medbf\char\appno.\the\sappno. #1}
	\nobreak\smallskip\noindent
}
%
% This is a subsection macro within the appendix.
%
\outer\def\ssapp #1\par{
	\advance\ssappno by 1
	\immediate\write\cntnts{\string\indent\string\indent\string\ninerm\ \appltr.\the\sappno.\the\ssappno\ #1\string\dotfill\folio\string\tenrm}
	\message{\appltr.\the\sappno.\the\ssappno}
	% Now write the section number onto the page:
	\vskip 0pt plus.1\vsize\penalty-250
	\smallskip
	\vskip 0pt plus-0.1\vsize
	\noindent
	{\bf\char\appno.\the\sappno.\the\ssappno. #1}
	\nobreak\smallskip\noindent
}
%
% Each time a reference is made, use the \ref macro. Calls the \references
% macro the first time that it is called.
%
\newif\iffirstref
\global\firstreftrue
\newcount\refno		\global\refno=0
%
%
% A macro to setup for references 
%
\def\references {
	\vfill\supereject
	\xdef\chapname{\llap{R}eferences}
	\immediate\write\cntnts{\string\smallskip}
	\immediate\write\cntnts{\string\noindent\string\bf\ References\string\tenrm\string\dotfill\folio}
	\frenchspacing
	\titlepage
	\noindent
	{\chapnumfont R}{\bigbf eferences}
	\bigskip
}
%
%
\def\ref #1 \par{
	\advance\refno by 1
	\iffirstref\references\firstreffalse\fi
	\item{[\the\refno]} #1
}
%
% Included here is a simple macro for putting figures into the text.
%
\def\figure #1 #2 #3 {
	\goodbreak\midinsert
	\vskip #2 cm
	\centerline{Figure #1.\ #3.}
	\endinsert
}
%
%
% Macros for making flow diagrams.
%
\def\mapright#1{\smash{
	\mathop{\longrightarrow}\limits^{#1}}}
\def\mapdown#1{\Big\downarrow
	\rlap{$\vcenter{\hbox{$\scriptstyle#1$}}$}}
\def\mapup#1{\Big\uparrow
	\rlap{$\vcenter{\hbox{$\scriptstyle#1$}}$}}
%
% This macro will copy a file EXACTLY as it is
%
\def\uncatcodespecials{\def\do##1{\catcode`##1=12 }\dospecials}
%
\def\setupverbatim{\ninett \lineno=0
 \def\par{\leavevmode\egroup\box0\endgraf}
 \obeylines \uncatcodespecials \obeyspaces
 \catcode`\`=\active \catcode`\^^I=\active
 \everypar{\advance\lineno by1
  \llap{\sevenrm\the\lineno\ \ }\startbox}}
{\catcode`\^^I=\active
 \gdef^^I{\leavevmode\egroup
  \dimen0=\wd0 % width so far, or since previous tab
  \divide\dimen0 by\w
  \multiply\dimen0 by\w %compute previous multiple of \w
  \advance\dimen0 by\w  %advance to next multiple of \w
  \wd0=\dimen0 \box0 \startbox}}
{\obeyspaces\global\let =\ }
\newdimen\w \setbox0=\hbox{\ninett\space} \w=8\wd0 %tab amount
\def\startbox{\setbox0=\hbox\bgroup}
%
\def\listing#1 {\par\begingroup\openup-1\jot\parindent=12pt
	\parskip=1.3pt plus 0.8pt
	\setupverbatim\input#1 \openup1\jot\parindent=20pt
	\parskip=2.0pt plus 1.0pt\endgroup}
%
